% !TEX root = ../main.tex
%% Corrigé du devoir surveillé n° 5 — valeurs vérifiées : verifications/devoirs.py

\corrigeDevoir{5}

\begin{corrigeBac}{7}{Le vélo à assistance électrique}
\begin{questions}
\item $E=U\,Q=36\times10=\res{\qty{360}{W.h}}$, soit $360\times3600=\res{\qty{1,30e6}{J}}$. \pts{1 pt}
\item $I=\dfrac{P}{U}=\dfrac{250}{36}=\res{\qty{6,9}{A}}$. \pts{1 pt}
\item $U=E-r\,I=37{,}0-0{,}15\times6{,}9=\res{\qty{36,0}{V}}$ : \res{\text{vérifié}}. Puissance dissipée : $P_\mathrm{J}=r\,I^2=0{,}15\times6{,}9^2=\res{\qty{7,1}{W}}$. \pts{1,5 pt}
\item $P_\mathrm{méca}=\eta\,P_\mathrm{reçue}=0{,}80\times250=\res{\qty{200}{W}}$ ; les \qty{50}{W} restants sont \res{\text{dissipés sous forme thermique}}
(effet Joule dans les bobinages, frottements). \pts{1 pt}
\item $\Delta t=\dfrac{E}{P}=\dfrac{360}{250}=\res{\qty{1,44}{h}}$ ; $d=25\times1{,}44=\res{\qty{36}{km}}$. \pts{1,5 pt}
\item $\qty{0,360}{kW.h}\times0{,}25=\res{0{,}09}$~€, moins de dix centimes. \pts{1 pt}
\end{questions}
\retiens{Une capacité en A.h multipliée par une tension en V donne une énergie en W.h ; $\qty{1}{W.h}=\qty{3600}{J}$.}
\end{corrigeBac}

\begin{corrigeBac}{7}{La descente à ski}
\begin{questions}
\item $h=\mathrm{AB}\,\sin\ang{15}=200\times\sin\ang{15}=\res{\qty{51,8}{m}}$. \pts{0,5 pt}
\item $W_\mathrm{AB}(\vec P)=m\,g\,(z_\mathrm{A}-z_\mathrm{B})=70\times9{,}81\times51{,}8=\res{\qty{3,55e4}{J}}$, travail moteur. \pts{1 pt}
\item Théorème de l'énergie cinétique, la réaction normale ne travaillant pas : $\dfrac12\,m\,v_\mathrm{B}^2-0=m\,g\,h$, d'où
$v_\mathrm{B}=\sqrt{2\,g\,h}=\res{\qty{31,9}{m.s^{-1}}}$, soit \res{\qty{115}{km/h}}. \pts{1,5 pt}
\item $E_\mathrm{c}(\mathrm{B})=\dfrac12\times70\times20{,}0^2=\res{\qty{1,40e4}{J}}$. Théorème : $E_\mathrm{c}(\mathrm{B})-0=W_\mathrm{AB}(\vec P)+W_\mathrm{AB}(\vec f)$, d'où
$W_\mathrm{AB}(\vec f)=\num{1,40e4}-\num{3,55e4}=\res{\qty{-2,15e4}{J}}$. \pts{1,5 pt}
\item $W_\mathrm{AB}(\vec f)=-f\times\mathrm{AB}$, donc $f=\dfrac{\num{2,15e4}}{200}=\res{\qty{108}{N}}$. \pts{1 pt}
\item $\Delta E_\mathrm{m}=W_\mathrm{AB}(\vec f)=\qty{-2,15e4}{J}$ : le skieur perd $\dfrac{\num{2,15e4}}{\num{3,55e4}}\approx\res{\qty{61}{\percent}}$ de son énergie
mécanique initiale (prise nulle en B pour l'énergie potentielle), transférée sous forme \res{\text{thermique}} à la neige, aux skis et
à l'air. \pts{1,5 pt}
\end{questions}
\schema[corrige]{ds5-piste}
\end{corrigeBac}

\begin{corrigeBac}{6}{La balle qui rebondit}
\begin{questions}
\item Elle calcule la vitesse au point $i$ à partir des positions qui l'encadrent : $v_i=\dfrac{y_{i-1}-y_{i+1}}{2\,\Delta t}$. La première et la dernière
position n'ont pas de voisine des deux côtés : on ne peut pas y calculer la vitesse. \pts{1,5 pt}
\item L'énergie cinétique, l'énergie potentielle de pesanteur (origine au sol) et l'énergie mécanique, en joules. Deuxième
position : $E_\mathrm{p}=m\,g\,y=0{,}058\times9{,}81\times1{,}488=\res{\qty{0,847}{J}}$. \pts{1 pt}
\item $E_\mathrm{m}$ reste constante (\qty{0,853}{J} à \qty{0,001}{J} près) : l'énergie potentielle se convertit en énergie cinétique ; les
\res{\text{frottements de l'air sont négligeables}} pendant la chute. \pts{1 pt}
\item Au sommet, la vitesse est nulle : $E_\mathrm{m}=m\,g\,y$. Avant : $0{,}058\times9{,}81\times1{,}500=\qty{0,853}{J}$ ; après : $0{,}058\times9{,}81\times1{,}10=\qty{0,626}{J}$.
Énergie perdue : \res{\qty{0,228}{J}}, soit $\dfrac{0{,}40}{1{,}50}\approx\res{\qty{27}{\percent}}$. \pts{1,5 pt}
\item Lors du choc, la balle et le sol se déforment : l'énergie est transférée sous forme \res{\text{thermique}} (échauffement) et
\res{\text{sonore}} (bruit du rebond). \pts{1 pt}
\end{questions}
\end{corrigeBac}
